\chapter{Basic Antiderivatives and Linearity}

The indefinite integral $\int f(x)\,\mathrm{d}x$ asks: which functions $F$
have derivative $f$? Such an $F$ is called an antiderivative of $f$, and any two antiderivatives differ by a
constant. Mastering this section means recognising at sight the
antiderivative of every standard building block (powers, $1/x$,
$\mathrm{e}^{ax}$) and using linearity to break compound integrands
apart. In economics, integration recovers a total quantity (total cost,
total revenue, total utility) from its marginal counterpart.

\section{Core Rules}

\begin{rulebox}{Definition}
A function $F$ is an antiderivative of $f$ on an interval if $F'(x)=f(x)$
there. The indefinite integral is the family of all antiderivatives:
\[
  \int f(x)\,\mathrm{d}x \;=\; F(x)+C,\qquad C\in\R .
\]
Always include the integration constant $C$ when writing an indefinite
integral.
\end{rulebox}

\begin{rulebox}{Power rule (reverse)}
For $n\in\R$ with $n\neq -1$,
\[
  \int x^{n}\,\mathrm{d}x \;=\; \frac{x^{n+1}}{n+1}+C .
\]
The case $n=-1$ is special and uses the logarithm.
\end{rulebox}

\begin{rulebox}{Logarithm}
For $x\neq 0$,
\[
  \int \frac{1}{x}\,\mathrm{d}x \;=\; \ln|x|+C .
\]
More generally, for constants $a\neq 0$ and $b$,
\[
  \int \frac{1}{ax+b}\,\mathrm{d}x \;=\; \frac{1}{a}\ln|ax+b|+C .
\]
\end{rulebox}

\begin{rulebox}{Exponential}
For any real $a\neq 0$,
\[
  \int \mathrm{e}^{x}\,\mathrm{d}x=\mathrm{e}^{x}+C,\qquad
  \int \mathrm{e}^{ax}\,\mathrm{d}x=\frac{1}{a}\mathrm{e}^{ax}+C .
\]
For a general base $b>0$, $b\neq 1$:
$\int b^{x}\,\mathrm{d}x=\dfrac{b^{x}}{\ln b}+C$.
\end{rulebox}

\begin{rulebox}{Linearity}
For constants $\alpha,\beta\in\R$ and integrable $f,g$:
\begin{itemize}[leftmargin=1.6em,itemsep=0.3ex]
  \item \textbf{Scalar:} $\displaystyle\int \alpha f(x)\,\mathrm{d}x
        =\alpha\int f(x)\,\mathrm{d}x$
  \item \textbf{Sum:} $\displaystyle\int\bigl(\alpha f(x)+\beta g(x)\bigr)\mathrm{d}x
        =\alpha\int f(x)\,\mathrm{d}x+\beta\int g(x)\,\mathrm{d}x$
\end{itemize}
There is \emph{no} product or quotient rule for integrals; these structures
require substitution, parts, or partial fractions (next sections).
\end{rulebox}

\begin{rulebox}{Useful rewrites}
Before integrating, rewrite the integrand:
\begin{itemize}[leftmargin=1.6em,itemsep=0.3ex]
  \item radicals as fractional powers: $\sqrt{x}=x^{1/2}$,
        $\sqrt[3]{x^{2}}=x^{2/3}$;
  \item reciprocals as negative powers: $1/x^{3}=x^{-3}$;
  \item split a sum over a single denominator into separate terms.
\end{itemize}
\end{rulebox}

\section{Practice Problems}

Compute each indefinite integral; include the constant of integration $C$.

\begin{problems}
  \item $\displaystyle\int x^{5}\,\mathrm{d}x$
  \item $\displaystyle\int 3x^{2}\,\mathrm{d}x$
  \item $\displaystyle\int \bigl(4x^{3}-2x+1\bigr)\mathrm{d}x$
  \item $\displaystyle\int \sqrt{x}\,\mathrm{d}x$
  \item $\displaystyle\int \frac{1}{x^{2}}\,\mathrm{d}x$
  \item $\displaystyle\int \frac{1}{x}\,\mathrm{d}x$
  \item $\displaystyle\int \mathrm{e}^{x}\,\mathrm{d}x$
  \item $\displaystyle\int \mathrm{e}^{2x}\,\mathrm{d}x$
  \item $\displaystyle\int \mathrm{e}^{-rt}\,\mathrm{d}t$
        \quad ($r$ constant, $r\neq 0$)
  \item $\displaystyle\int \frac{1}{2x+3}\,\mathrm{d}x$
  \item $\displaystyle\int \bigl(5\mathrm{e}^{x}-3x^{2}\bigr)\mathrm{d}x$
  \item $\displaystyle\int \frac{1}{\sqrt{x}}\,\mathrm{d}x$
  \item $\displaystyle\int \frac{x^{2}+1}{x}\,\mathrm{d}x$
  \item $\displaystyle\int \bigl(x^{3}+2\bigr)^{2}\mathrm{d}x$
        \quad \emph{(expand first)}
  \item $\displaystyle\int 7\,\mathrm{d}x$
  \item $\displaystyle\int \sqrt[3]{x}\,\mathrm{d}x$
  \item $\displaystyle\int \frac{2}{x}\,\mathrm{d}x$
  \item $\displaystyle\int \bigl(6x^{5}-4x^{3}+x\bigr)\mathrm{d}x$
  \item Marginal cost: $C'(q)=4+0.1q$. Find $C(q)$ given $C(0)=100$.
  \item $\displaystyle\int \frac{x^{3}-2x+5}{x^{2}}\,\mathrm{d}x$
  \item $\displaystyle\int \Bigl(\sqrt{x}+\frac{1}{\sqrt{x}}\Bigr)\mathrm{d}x$
  \item $\displaystyle\int \mathrm{e}^{-x}\,\mathrm{d}x$
  \item $\displaystyle\int 5\mathrm{e}^{3x}\,\mathrm{d}x$
  \item $\displaystyle\int \frac{3}{4x-1}\,\mathrm{d}x$
  \item Marginal revenue: $R'(q)=20-0.2q$. Find $R(q)$ given $R(0)=0$.
\end{problems}

\section{Solutions}

\begin{solutions}
  \item Power rule with $n=5$:
        $\displaystyle\int x^{5}\,\mathrm{d}x=\frac{x^{6}}{6}+C$.

  \item $\displaystyle\int 3x^{2}\,\mathrm{d}x
        =3\cdot\frac{x^{3}}{3}+C=x^{3}+C$.

  \item Term by term:
        $\displaystyle\int\bigl(4x^{3}-2x+1\bigr)\mathrm{d}x
         =x^{4}-x^{2}+x+C$.

  \item $\sqrt{x}=x^{1/2}$, so
        \[
          \int x^{1/2}\,\mathrm{d}x=\frac{x^{3/2}}{3/2}+C
            =\tfrac{2}{3}x^{3/2}+C .
        \]

  \item $1/x^{2}=x^{-2}$, so
        $\displaystyle\int x^{-2}\,\mathrm{d}x=\frac{x^{-1}}{-1}+C
         =-\frac{1}{x}+C$.

  \item $\displaystyle\int \frac{1}{x}\,\mathrm{d}x=\ln|x|+C$.

  \item $\displaystyle\int \mathrm{e}^{x}\,\mathrm{d}x=\mathrm{e}^{x}+C$.

  \item $\displaystyle\int \mathrm{e}^{2x}\,\mathrm{d}x
        =\tfrac12\mathrm{e}^{2x}+C$.

  \item $\displaystyle\int \mathrm{e}^{-rt}\,\mathrm{d}t
        =-\frac{1}{r}\mathrm{e}^{-rt}+C$.

  \item $\displaystyle\int \frac{1}{2x+3}\,\mathrm{d}x
        =\tfrac12\ln|2x+3|+C$.

  \item Linearity:
        \[
          \int\bigl(5\mathrm{e}^{x}-3x^{2}\bigr)\mathrm{d}x
            =5\mathrm{e}^{x}-x^{3}+C .
        \]

  \item $1/\sqrt{x}=x^{-1/2}$, so
        \[
          \int x^{-1/2}\,\mathrm{d}x=\frac{x^{1/2}}{1/2}+C=2\sqrt{x}+C .
        \]

  \item Split: $\dfrac{x^{2}+1}{x}=x+\dfrac{1}{x}$, so
        \[
          \int\Bigl(x+\frac{1}{x}\Bigr)\mathrm{d}x
            =\frac{x^{2}}{2}+\ln|x|+C .
        \]

  \item Expand: $(x^{3}+2)^{2}=x^{6}+4x^{3}+4$, so
        \[
          \int\bigl(x^{6}+4x^{3}+4\bigr)\mathrm{d}x
            =\frac{x^{7}}{7}+x^{4}+4x+C .
        \]

  \item $\displaystyle\int 7\,\mathrm{d}x=7x+C$.

  \item $\sqrt[3]{x}=x^{1/3}$, so
        $\displaystyle\int x^{1/3}\,\mathrm{d}x=\frac{x^{4/3}}{4/3}+C
         =\tfrac{3}{4}x^{4/3}+C$.

  \item $\displaystyle\int \frac{2}{x}\,\mathrm{d}x=2\ln|x|+C$.

  \item Term by term:
        $\displaystyle\int\bigl(6x^{5}-4x^{3}+x\bigr)\mathrm{d}x
         =x^{6}-x^{4}+\tfrac12 x^{2}+C$.

  \item Antidifferentiate $C'(q)$:
        \[
          C(q)=\int\bigl(4+0.1q\bigr)\mathrm{d}q=4q+0.05q^{2}+K .
        \]
        The condition $C(0)=100$ gives $K=100$, so
        $C(q)=100+4q+0.05q^{2}$.

  \item Split the integrand:
        $\dfrac{x^{3}-2x+5}{x^{2}}=x-\dfrac{2}{x}+\dfrac{5}{x^{2}}$. Then
        \[
          \int\Bigl(x-\frac{2}{x}+5x^{-2}\Bigr)\mathrm{d}x
            =\frac{x^{2}}{2}-2\ln|x|-\frac{5}{x}+C .
        \]

  \item $\displaystyle\int\bigl(x^{1/2}+x^{-1/2}\bigr)\mathrm{d}x
        =\tfrac{2}{3}x^{3/2}+2x^{1/2}+C$.

  \item $\displaystyle\int \mathrm{e}^{-x}\,\mathrm{d}x=-\mathrm{e}^{-x}+C$.

  \item $\displaystyle\int 5\mathrm{e}^{3x}\,\mathrm{d}x
        =\tfrac{5}{3}\mathrm{e}^{3x}+C$.

  \item $\displaystyle\int \frac{3}{4x-1}\,\mathrm{d}x
        =\tfrac{3}{4}\ln|4x-1|+C$.

  \item Antidifferentiate:
        \[
          R(q)=\int\bigl(20-0.2q\bigr)\mathrm{d}q=20q-0.1q^{2}+K .
        \]
        With $R(0)=0$ we get $K=0$, so $R(q)=20q-0.1q^{2}$ (a standard
        linear-demand revenue).
\end{solutions}

\section{Additional Problems}

Compute each indefinite integral; include the constant of integration $C$.

\begin{problems}
  \item $\displaystyle\int x^{8}\,\mathrm{d}x$
  \item $\displaystyle\int 10x^{4}\,\mathrm{d}x$
  \item $\displaystyle\int \bigl(2x^{5}-3x^{2}+7\bigr)\mathrm{d}x$
  \item $\displaystyle\int \sqrt[4]{x}\,\mathrm{d}x$
  \item $\displaystyle\int \frac{4}{x^{3}}\,\mathrm{d}x$
  \item $\displaystyle\int \frac{3}{x}\,\mathrm{d}x$
  \item $\displaystyle\int \mathrm{e}^{-2x}\,\mathrm{d}x$
  \item $\displaystyle\int \mathrm{e}^{0.04t}\,\mathrm{d}t$
  \item $\displaystyle\int \frac{1}{5x-2}\,\mathrm{d}x$
  \item $\displaystyle\int \bigl(2\mathrm{e}^{x}+x^{4}\bigr)\mathrm{d}x$
  \item $\displaystyle\int \frac{x^{4}-3x^{2}}{x^{2}}\,\mathrm{d}x$
  \item $\displaystyle\int \bigl(x+1\bigr)^{2}\mathrm{d}x$
        \quad \emph{(expand first)}
  \item $\displaystyle\int \Bigl(\sqrt{x}-\frac{2}{x}\Bigr)\mathrm{d}x$
  \item Marginal cost: $C'(q)=6+0.4q$. Find $C(q)$ given $C(0)=80$.
  \item Marginal revenue: $R'(q)=50-0.5q$. Find $R(q)$ given $R(0)=0$.
\end{problems}

\section{Additional Solutions}

\begin{solutions}
  \item $\displaystyle\int x^{8}\,\mathrm{d}x=\frac{x^{9}}{9}+C$.

  \item $\displaystyle\int 10x^{4}\,\mathrm{d}x=2x^{5}+C$.

  \item Term by term:
        \[
          \int\bigl(2x^{5}-3x^{2}+7\bigr)\mathrm{d}x
            =\frac{x^{6}}{3}-x^{3}+7x+C .
        \]

  \item $\sqrt[4]{x}=x^{1/4}$, so
        \[
          \int x^{1/4}\,\mathrm{d}x=\frac{x^{5/4}}{5/4}+C
            =\tfrac{4}{5}x^{5/4}+C .
        \]

  \item $4/x^{3}=4x^{-3}$, so
        $\displaystyle\int 4x^{-3}\,\mathrm{d}x
         =4\cdot\frac{x^{-2}}{-2}+C=-\frac{2}{x^{2}}+C$.

  \item $\displaystyle\int \frac{3}{x}\,\mathrm{d}x=3\ln|x|+C$.

  \item $\displaystyle\int \mathrm{e}^{-2x}\,\mathrm{d}x
        =-\tfrac12\mathrm{e}^{-2x}+C$.

  \item $\displaystyle\int \mathrm{e}^{0.04t}\,\mathrm{d}t
        =\frac{1}{0.04}\mathrm{e}^{0.04t}+C=25\,\mathrm{e}^{0.04t}+C$.

  \item $\displaystyle\int \frac{1}{5x-2}\,\mathrm{d}x
        =\tfrac{1}{5}\ln|5x-2|+C$.

  \item Linearity:
        \[
          \int\bigl(2\mathrm{e}^{x}+x^{4}\bigr)\mathrm{d}x
            =2\mathrm{e}^{x}+\tfrac{x^{5}}{5}+C .
        \]

  \item Split: $\dfrac{x^{4}-3x^{2}}{x^{2}}=x^{2}-3$, so
        \[
          \int\bigl(x^{2}-3\bigr)\mathrm{d}x=\tfrac{x^{3}}{3}-3x+C .
        \]

  \item Expand: $(x+1)^{2}=x^{2}+2x+1$, so
        \[
          \int\bigl(x^{2}+2x+1\bigr)\mathrm{d}x
            =\tfrac{x^{3}}{3}+x^{2}+x+C .
        \]

  \item $\displaystyle\int\bigl(x^{1/2}-2/x\bigr)\mathrm{d}x
        =\tfrac{2}{3}x^{3/2}-2\ln|x|+C$.

  \item Antidifferentiate $C'(q)$:
        \[
          C(q)=\int\bigl(6+0.4q\bigr)\mathrm{d}q=6q+0.2q^{2}+K .
        \]
        With $C(0)=80$, $K=80$, so $C(q)=80+6q+0.2q^{2}$.

  \item Antidifferentiate:
        \[
          R(q)=\int\bigl(50-0.5q\bigr)\mathrm{d}q=50q-0.25q^{2}+K .
        \]
        With $R(0)=0$ we get $K=0$, so $R(q)=50q-0.25q^{2}$.
\end{solutions}
